\section{Heuristic A*}
\subsection{Định nghĩa heuristic}
Với mỗi scenario \(s\), nhóm em tính \(\rho_s\) một lần trước khi bắt đầu tìm kiếm:
\[
\rho_s=\min_{\substack{e\in E,\ e\text{ không đóng}\\
\operatorname{Haversine}(e)>10^{-6}\ {\rm km}}}
\frac{\operatorname{Cost}_s(e)}{\operatorname{Haversine}(e)}.
\]
Sau đó:
\[
h_s(n)=\rho_s\cdot\operatorname{Haversine}(n,\mathrm{goal}),
\qquad f(n)=g(n)+h_s(n).
\]
Edge đã đóng không được đưa vào phép tính \(\rho_s\). Edge có khoảng cách Haversine gần 0 cũng được bỏ qua để tránh chia cho số quá nhỏ. Nếu không còn edge hợp lệ thì \(\rho_s=0\); node thiếu tọa độ có \(h(n)=0\).

\subsection{Admissibility và consistency}
Vì \(\rho_s\) không lớn hơn tỉ số giữa cost và khoảng cách của mọi edge hợp lệ, đồng thời Haversine thỏa bất đẳng thức tam giác, ta có:
\[
h_s(n)\leq \operatorname{cost}(\text{một đường đi từ }n\text{ đến goal}).
\]
Do đó \(h_s(n)\geq0\) và heuristic là admissible. Với một edge \(e=(u,v)\), ta cũng có:
\[
h_s(u)\leq \operatorname{Cost}_s(e)+h_s(v),
\]
nên heuristic là consistent. Khi edge cost không âm, A* trả về route có composite cost tối ưu, giống như UCS.

\subsection{Ảnh hưởng của scenario}
Scenario được giữ cố định trong một lần tìm kiếm nên cost edge và \(\rho_s\) không thay đổi giữa các bước mở rộng node. Khi traffic, flood hoặc weight thay đổi, cần tính lại \(\rho_s\). Heuristic sử dụng cost của edge sau khi đã kết hợp distance, free-flow time, congestion và flood risk; do đó không bị dư dữ kiện mà đang tạo lower bound cho đúng composite cost của scenario. Heuristic này không tự biến bài toán thành tối ưu khoảng cách thuần túy: A* tối ưu composite cost, còn mục tiêu ưu tiên đường ngắn nhất được thể hiện bằng cost cơ sở và cách diễn giải kết quả.
